\documentclass[]{article}

\usepackage[T1]{fontenc}
\usepackage[utf8]{inputenc}
\usepackage{amsthm, amsmath, amsfonts}
\usepackage{blkarray}
\usepackage{bigstrut}
\usepackage{graphicx}

\newcommand{\avaliacao}{3}
\newcommand{\disciplina}{Probabilidade e Estatística}
\newcommand{\semestre}{2024/1}


\theoremstyle{definition}
\newtheorem{ex}{Exercício}

\renewcommand{\SS}{\mathcal{S}}

\begin{document}
\thispagestyle{empty}
\begin{center}
\Large{\avaliacao ª Avaliação de \disciplina\ - \semestre}

\large{Prof. Dr. Alessandro José Queiroz Sarnaglia}
\end{center}

\begin{ex}
A resistência à ruptura de rebites de um fabricante possui valor médio de 9000 psi e desvio padrão de 400 psi.
\begin{enumerate}
	\item Qual é a probabilidade de a resistência à ruptura média de uma amostra aleatória de 36 rebites estar entre 8900 e 9100?
	\item Se o tamanho da amostra fosse 81, e não 36, qual seria a probabilidade do item anterior?
\end{enumerate}
\end{ex}


\begin{ex}
Seja $X_1, X_2, \ldots, X_n$ uma amostra aleatória de uma distribuição de Gama($\alpha$, $\beta$), com $\alpha=3$. Isto é, a amostra é coletada de uma população com fdp
$$
f(x;\theta) = \frac{1}{\beta^3\Gamma(3)}x^{2}\exp\left\{-\frac{x}{\beta}\right\}, \ \ x > 0, \ \ \beta>0.
$$
\begin{enumerate}
	\item Encontre $\hat\beta$, o Estimador de Máxima Verossimilhança de $\beta$.
	\item Qual é o vício de $\hat\beta$? \underline{Dica: $E(\bar{X}) = \mu$.}
	\item Estime $\beta$ a partir das $n = 10$ observações a seguir sobre a força vibratória de uma palheta de turbina sob condições especificadas:
	\begin{center}
	\begin{tabular}{ccccc}
		27,8 & 10,6 & 14,4 & 15,3 & 21,2 \\
		6,5 & 15,5 & 23,6 & 29,2 & 7,9\\
	\end{tabular}
	\end{center}
\end{enumerate}
\end{ex}


\begin{ex}
Um fabricante de aspirina enche os frascos por peso em vez de fazê-lo por quantidade. Uma vez que cada frasco deve conter 100 comprimidos, o peso médio por unidade deve ser de 5 \textit{grains}. Cada um dos 100 comprimidos tirados de um lote bastante grande é pesado, resultando em um peso médio amostral por comprimido de 4,87 \textit{grains} e uma variância da amostra de 0,49 \textit{grains}.
\begin{enumerate}
	\item Calcule o $p$-valor para o teste das hipóteses $H_0 : \mu \geq 5 \times H_1 : \mu < 5$.
	\item Utilizando o $p$-valor obtido acima e um nível de significância de 0,01, você rejeitaria $H_0$ ou não? Justifique.
\end{enumerate}
\end{ex}


\end{document}
